% Direct-authority diplomatic French for printed p.181.
\SourcePage{181}

\LayerTarget{tohoku:proof:proposition-3.9.2:finite-intersection:p2}
des $E_x(\mathbf R)$ pour $\mathbf R$ parcourant $\mathfrak R(X)$ est par
hypothèse réduite à $x$, d'oû on conclut, par raison de quasi-compacité,
qu'il existe une partie finie $L_x$ de $\mathfrak R(X)$ telle que
l'intersection des $E_x(\mathbf R)$ pour $\mathbf R\in L_x$, i.e.
$E_x(\mathbf R^{L_x})$, est contenue dans $U_x$. Les
$O_x(\mathbf R^{L_x})$ forment un recouvrement ouvert de $X$, il existe donc
un ensemble fini $Y\subset X$ tel que les $O_x(\mathbf R^{L_x})$
correspondants aux $x\in Y$ recouvrent $X$. Soit $L$ la réunion des $L_x$
pour $x\in Y$, je dis que le recouvrement $\mathbf R^L$ de $X$ est plus fin
que $\mathbf U$. Soit en effet $A\in\mathbf R^L$ non vide, soit $a\in A$,
il existe $x\in Y$ tel que $a\in O_x(\mathbf R^{L_x})$ d'où
$E_a(\mathbf R^{L_x})\subset E_x(\mathbf R^{L_x})\subset U_x$ et à fortiori
$A\subset E_a(\mathbf R^L)\subset E_a(\mathbf R^{L_x})$ est contenu dans
$U_x$, ce qui démontre notre assertion. Appliquant ce résultat, pour tout
$A\in\mathfrak S$, à l'ensemble des recouvrements de $A$ induits par les
recouvrements $\mathbf R\in\mathfrak R(X)$, la conclusion voulue apparaît.

\LayerTarget{tohoku:application:corollary-3.9.1-3:cubes}
Le cas d'application le plus frappant du corollaire 3 est celui où $X$ est
le cube compact $0\leqq x_i\leqq1$ de $\mathbf R^n$, oû $\mathfrak S$ est la
famille des cubes compacts $A$ du type $a_i\leqq x_i\leqq b_i$ contenus dans
$X$, $\mathfrak R(A)$ étant la famille des recouvrements à ensembles deux
$A$ définis par des hyperplans parallèles aux hyperplans coordonnées : Pour
vérifier que $H^i(A,F)=0$ pour $i>0$ et tout $A$, il suffit de vérifier que
pour tout $A$, et toute section $f$ de $F$ sur $A_1\cap A_2$, on a
$f=f_1-f_2$, oû $f_i$ est une section de $F$ sur $A_i$. C'est la réduction
faite par H. Cartan dans sa démonstration des théorèmes fondamentaux sur les
variétés de Stein [5].

\LayerTarget{tohoku:remark:theorem-3.9.1:nonabelian-n-one}
\textsc{Remarque.} Si $n=1$, le théorème 3.9.1 garde un sens, et se
vérifie facilement directement, si on suppose que $F$ est un faisceau de
groupes non nécessairement abéliens. Cela permet de simplifier la
démonstration du théorème [5, XVIII] sur les matrices holomorphes inversibles.

\LayerTarget{tohoku:section:3.10}
\textbf{3.10. Passages à la limite en cohomologie des faisceaux.}
Nous ne donnerons que deux résultats dans cette voie, (dont l'un nous servira
au Chapitre 5, N\textsuperscript{o} 7), cas particuliers du résultat général
suivant d'algèbre homologique :

\LayerTarget{tohoku:proposition:3.10.1}
\textsc{Proposition 3.10.1.} \textit{Soient $\mathbf C,\mathbf C'$ deux
catégories abéliennes, on suppose que tout élément de $\mathbf C$ est
isomorphe à un sous-truc d'un élément injectif, et que $\mathbf C'$ satisfait
à l'axiome \textup{AB 5 (}cf. \textup{1.5)}, qui permet en particulier de prendre des limites
inductives dans $\mathbf C'$ \textup{(}cf. prop. \textup{1.8)}. Soit $(F_i)_{i\in I}$ un
système inductif de foncteurs additifs covariants de $\mathbf C$ dans
$\mathbf C'$, soit $F=\mathop{\underset{\to}{\lim}} F_i$ le foncteur limite inductive des $F_i$,
défini par $F(A)=\mathop{\underset{\to}{\lim}} F_i(A)$ pour tout $A\in\mathbf C$. Les
homomorphismes $F_i\to F$ définissent des homomorphismes de
$\partial$-foncteurs $(\mathbf R^pF_i)\to(\mathbf R^pF)$, d'oû
un homomorphisme de $\partial$-foncteurs :}
\LayerTarget{tohoku:eq:3.10.1}
\[
\tag{3.10.1.}
\mathop{\underset{\to}{\lim}} \mathbf R^pF_i(A)\to \mathbf R^pF(A)
\]
\textit{\textup{(}les homomorphismes cobord pour la suite de foncteurs
$\mathop{\underset{\to}{\lim}}\mathbf R^pF_i$ se définissant comme limite inductive des
homomorphismes cobords relatifs aux $\mathbf R^pF_i$\textup{)}. Les homomorphismes
\textup{(3.10.1)} sont des isomorphismes.}

\LayerTarget{tohoku:proof:proposition-3.10.1:p1}
Il suffit pour le voir de prendre une résolution injective
$C=C(A)$ de $A$, le premier membre de (3.10.1) est alors
$\mathop{\underset{\to}{\lim}} H^p(F_iC(A))$, le deuxième est

\TohokuContinuation{182}
